% from ega4-2-fr.tex
\oldpage[IV-2]{7}
$Y'=\operatorname{Spec}(\mathcal O_{y'})$ et
$Z=\operatorname{Spec}(\mathcal O_z)$, et il suffit alors de prouver que
$\mathscr F\otimes_{\mathcal O_{Y'}}\mathscr G'$ est $(h\circ f')$-plat. Compte tenu de
(2.1.2), la proposition résulte de Bourbaki, \emph{Alg. comm.},
chap.~I\textsuperscript{er}, \S~2, n\textsuperscript{o}~7, prop.~8.

\begin{corollary}[2.1.6]
\label{IV.2.1.6-fr}
Soient $f:X\to Y$, $g:Y\to Z$ deux morphismes de préschémas, $\mathscr F$ un
$\mathcal O_X$-Module. Si $\mathscr F$ est $f$-plat au point $x\in X$ et si
$g$ est plat au point $f(x)$, $\mathscr F$ est $(g\circ f)$-plat au point $x$. En
particulier, si $f$ et $g$ sont des morphismes plats, il en est de même de
$g\circ f$.
\end{corollary}

C'est le cas particulier de (2.1.5) avec $Y'=Y$, $\mathscr G=\mathcal O_{Y'}$.

\begin{corollary}[2.1.7]
\label{IV.2.1.7-fr}
Si $f:X\to X'$, $g:Y\to Y'$ sont deux $S$-morphismes plats,
$f\times_S g:X\times_S Y\to X'\times_S Y'$ est un morphisme plat.
\end{corollary}

Cela résulte de (2.1.4) et (2.1.6) (cf. \textbf{I}, 3.5.1).

\begin{proposition}[2.1.8]
\label{IV.2.1.8-fr}
Soient $f:X\to Y$ un morphisme de préschémas,
\[
0\to\mathscr F'\to\mathscr F\to\mathscr F''\to0
\]
une suite exacte de $\mathcal O_X$-Modules quasi-cohérents, telle que
$\mathscr F''$ soit $Y$-plat.
\begin{enumerate}
\item[\textup{(i)}] Pour tout morphisme $g:Y'\to Y$ et tout $\mathcal O_{Y'}$-Module
quasi-cohérent $\mathscr G'$, la suite
$0\to\mathscr F'\otimes_Y\mathscr G'\to
\mathscr F\otimes_Y\mathscr G'\to
\mathscr F''\otimes_Y\mathscr G'\to0$
de $\mathcal O_{X'}$-Modules (où $X'=X\times_Y Y'$) est exacte.
\item[\textup{(ii)}] Pour que $\mathscr F$ soit $Y$-plat, il faut et il suffit que
$\mathscr F'$ le soit.
\end{enumerate}
\end{proposition}

On peut évidemment supposer $X$, $Y$, $Y'$ affines~; la conclusion résulte
alors de (2.1.2) et de $(\mathbf{0}_{\mathrm I},6.1.2)$.

\begin{corollary}[2.1.9]
\label{IV.2.1.9-fr}
Soient $\mathscr L^\bullet$ un complexe de $\mathcal O_X$-Modules
quasi-cohérents, $i$ un indice tel que si l'on note
$d^i:\mathscr L^i\to\mathscr L^{i+1}$ l'opérateur de dérivation,
$\mathscr B^{i+1}(\mathscr L^\bullet)=\operatorname{Im}(d^i)$ et
$\mathscr Z'^{i+1}(\mathscr L^\bullet)=\operatorname{Coker}(d^i)$ soient
$Y$-plats. Alors, avec les notations de \textup{(2.1.8)}, l'homomorphisme canonique
\[
\mathscr H^i(\mathscr L^\bullet)\otimes_Y\mathscr G'
\longrightarrow
\mathscr H^i(\mathscr L^\bullet\otimes_Y\mathscr G')
\]
est bijectif.
\end{corollary}

Comme le produit tensoriel est exact à droite, on a
\[
\mathscr Z'^{i+1}(\mathscr L^\bullet)\otimes_Y\mathscr G'
=\operatorname{Coker}(d^i\otimes1)
=\mathscr Z'^{i+1}(\mathscr L^\bullet\otimes_Y\mathscr G')
\]
et
$\mathscr Z'^i(\mathscr L^\bullet)\otimes_Y\mathscr G'
=\mathscr Z'^i(\mathscr L^\bullet\otimes_Y\mathscr G')$. En outre, dans la
suite exacte
\[
0\to\mathscr B^{i+1}(\mathscr L^\bullet)\to\mathscr L^{i+1}
\to\mathscr Z'^{i+1}(\mathscr L^\bullet)\to0
\]
$\mathscr Z'^{i+1}(\mathscr L^\bullet)$ est $Y$-plat, donc il résulte de
(2.1.8, (i)) que l'on a la suite exacte
\[
0\to\mathscr B^{i+1}(\mathscr L^\bullet)\otimes_Y\mathscr G'
\to\mathscr L^{i+1}\otimes_Y\mathscr G'
\to\mathscr Z'^{i+1}(\mathscr L^\bullet\otimes_Y\mathscr G')\to0
\]
d'où
$\mathscr B^{i+1}(\mathscr L^\bullet)\otimes_Y\mathscr G'
=\operatorname{Im}(d^i\otimes1)
=\mathscr B^{i+1}(\mathscr L^\bullet\otimes_Y\mathscr G')$. Alors, comme
dans la suite exacte
\[
0\to\mathscr H^i(\mathscr L^\bullet)\to\mathscr Z'^i(\mathscr L^\bullet)
\to\mathscr B^{i+1}(\mathscr L^\bullet)\to0
\]
$\mathscr B^{i+1}(\mathscr L^\bullet)$ est $Y$-plat, il résulte de
(2.1.8, (i)) et de ce qui précède que l'on a la suite exacte
\[
0\to\mathscr H^i(\mathscr L^\bullet)\otimes_Y\mathscr G'
\to\mathscr Z'^i(\mathscr L^\bullet\otimes_Y\mathscr G')
\to\mathscr B^{i+1}(\mathscr L^\bullet\otimes_Y\mathscr G')\to0
\]
ce qui démontre le corollaire.

